% ============================================================
% Chapter 2: State of the Art
%   Section: Positioning of This Work
% Included from chapters/02-state-of-the-art/state-of-the-art.tex
% ============================================================

\section{Positioning of This Work}
\label{sec:positioning}

Table~\ref{tab:positioning} sets the work against the closest
benchmarks and tools on the properties the problem statement
(Section~\ref{sec:problem}) requires.

\begin{table}[ht]
\centering
\small
\begin{tabularx}{\linewidth}{@{}Xcccccc@{}}
\toprule
& \rotatebox{70}{Open-web evidence} & \rotatebox{70}{English and Spanish}
& \rotatebox{70}{Constructive news} & \rotatebox{70}{Local, no paid API}
& \rotatebox{70}{Cited, checkable verdict} & \rotatebox{70}{Finds the news} \\
\midrule
FEVER, SciFact & -- & -- & -- & n/a & -- & -- \\
MultiFC, X-Fact & snippets & X-Fact & -- & n/a & -- & -- \\
AVeriTeC & yes & -- & -- & n/a & yes & -- \\
ClaimBuster & -- & -- & -- & yes & -- & -- \\
Web-search LLM assistants & yes & yes & -- & -- & partly & -- \\
\textbf{This work} & yes & yes & yes & yes & yes & yes \\
\bottomrule
\end{tabularx}
\caption{The work against the closest benchmarks and tools. ``n/a'':
a dataset, not a system.}
\label{tab:positioning}
\end{table}

The contribution is therefore not a new model or a new benchmark
result. It is an end-to-end pipeline that reads constructive news in
two languages, finds and ranks what is worth publishing on the topics
asked for, and verifies its claims on the open web with local models
and verdicts that can be audited --- \emph{evaluated stage by stage,
against a person wherever a person takes the same decision}. The
pipeline is the artefact; the evaluation is the contribution. Its
custom validation set is, to the author's knowledge, the first labelled
set of claims drawn from constructive news, and its X-Fact pilot
documents what years between a claim and the search for its evidence
do to an open-retrieval score.
